\section{Related Work}
\label{sec:related_work}
% CN-SECTION: 文献综述按主题说明AIM、JSSP/RL调度和轨迹平滑的已有工作及不足。
% EN: Organize this section by research themes rather than paper-by-paper summaries. Each paragraph should end by explaining what limitation motivates this paper.
% CN: 本节建议按主题组织，而不是逐篇论文罗列。每个主题最后说明它和本文工作的差距。

\subsection{Autonomous Intersection Management}
% EN: Review signal-free AIM, reservation-based approaches, and centralized/decentralized coordination.
% CN: 回顾无信号灯路口管理、预约式方法、集中式和分布式协调方法。
Autonomous intersection management replaces fixed traffic signals with negotiated or optimized vehicle movements and has been widely discussed as a promising direction for connected autonomous traffic systems \cite{zhong2020autonomous}. Reservation-based approaches divide the intersection into spatial or temporal resources, while rule-based methods determine priority through arrival time, lane, or safety margins. These methods are practical baselines because they are interpretable and computationally light, but they may become conservative when traffic demand is high.

\subsection{Scheduling and Graph-Based Formulations}
% EN: Explain how conflict zones can be viewed as shared machines and vehicles as jobs. Discuss why JSSP is attractive for dense traffic.
% CN: 解释冲突区如何对应机器、车辆如何对应作业，并说明 JSSP 适合高密度交通的原因。
The job shop scheduling problem provides a natural abstraction for intersection coordination. Each vehicle can be treated as a job, each conflict zone as a machine, and each vehicle-zone traversal as an operation. Precedence constraints preserve the geometric order of zones along a path, while disjunctive constraints prevent two conflicting vehicles from occupying the same zone at the same time. This abstraction is attractive because it converts a traffic coordination problem into a structured combinatorial optimization problem.

\subsection{Reinforcement Learning for Combinatorial Scheduling}
% EN: Discuss RL policies that construct schedules step by step, including state, action, and reward design.
% CN: 讨论强化学习如何逐步构建调度序列，重点写状态、动作和奖励设计。
Reinforcement learning has been used to learn dispatching policies for combinatorial scheduling problems where handcrafted priority rules are difficult to tune. In a JSSP setting, an agent can observe the current partial schedule and select the next schedulable operation. For intersection management, this mechanism can be adapted to select which vehicle should enter a conflict zone next. The key challenge is to design an observation representation that captures safety, route order, and traffic density while remaining computationally efficient.

\begin{addedblock}
In the present framework, reinforcement learning is used to learn a constructive scheduling policy before deployment. During online intersection control, the trained policy parameters remain fixed, and the policy is queried repeatedly to complete the schedule for the active vehicle set. This distinction is important because online rollout has a clearer real-time role than continuous online training while vehicles are being controlled.
\end{addedblock}

\begin{addedtwoblock}
The proposed scheduler uses a relation-aware graph attention encoder because the operation graph contains typed structural relations: route-precedence edges, reverse route context, and same-zone disjunctive conflicts. This design is motivated by graph attention networks, relation-aware attention, relational graph neural networks, and learning-based JSSP dispatching on disjunctive graphs \cite{velickovic2018graphattention,shaw2018selfattention,schlichtkrull2018rgcn,chen2021rgat,zhang2020learningdispatch}.
\end{addedtwoblock}
\subsection{Trajectory Smoothing and Stop Avoidance}
% EN: Review velocity planning, energy-aware control, control barrier functions, and smooth crossing strategies.
% CN: 回顾速度规划、能耗优化、控制屏障函数和平滑通过路口的方法。
\begin{addedtodayblock}
Recent work shows that intersection efficiency depends not only on the discrete passing order but also on whether the assigned order can be executed by dynamically feasible vehicle trajectories. Graph-based CAV cooperation methods model crossing, diverging, converging, and reachability conflicts before selecting a passing order, which is close in spirit to representing conflict zones as shared scheduling resources \cite{chen2021conflictfree,chen2021lanechanging}. Priority-aware resequencing further shows that strict FCFS/FIFO decision orders can be relaxed through replanning while still passing target times to a lower-level optimal controller \cite{chalaki2021priority}. These studies motivate the scheduling layer in this paper, but they usually rely on handcrafted graph algorithms or analytic resequencing rather than an offline-trained graph policy.

Trajectory-level studies address the complementary question of how a vehicle should move after a desired crossing time, phase, or maneuver has been selected. Eco-driving and eco-approach controllers use reinforcement learning, Pontryagin-minimum-principle solutions, or hybrid rule-learning structures to reduce energy consumption, delay, and unnecessary acceleration near intersections \cite{bai2022hybrid,jiang2023ecodriving,lakshmanan2022cooperative,dong2023overtaking}. These papers support the use of stop count, acceleration variation, and energy proxy metrics in the evaluation, because stop-line waiting is not merely a delay phenomenon but also changes the smoothness and energy profile of the vehicle trajectory.

Safety-oriented trajectory planning adds constraints that the high-level schedule alone cannot guarantee. Predictive-control and spatial-domain formulations derive trajectories under rear-end, crossing, merging, diverging, and mixed-traffic uncertainty constraints \cite{mahbub2022safety,zhao2025spatial}. Learning-based AIM approaches can improve flexibility and travel time, but recent work also emphasizes that learned decisions need explicit safety evaluation or filtering, such as uncertainty-aware high-order control barrier functions \cite{luo2023realtime,yu2025uncertainty,cederle2024distributed,li2023dhal}. A broad CAV-control survey similarly identifies trajectory tracking, collaborative control, safety, comfort, and energy saving as coupled requirements rather than independent design goals \cite{liu2023systematic}.

The proposed framework follows this schedule-to-trajectory view. The high-level JSSP/RL scheduler determines the conflict-zone order and target entry times, while the low-level smoother treats those times as references to be tracked under speed, acceleration, jerk, and safety constraints. If the smoother cannot track the assigned reservations without violating collision-avoidance or comfort bounds, it reports infeasibility and triggers replanning or conservative stop-line fallback. This feedback interface is the main distinction from methods that optimize only the passing order or only the vehicle trajectory in isolation.
\end{addedtodayblock}

% EN: Before submission, ensure every cited claim has a verified source. Do not add placeholder citations.
% CN: 投稿前必须确认每个引用支撑相应论断，不要为了占位编造引用。




